% Chapter 4 -- The acceptance procedure (PUB-REF-0001 AN-15 procedure
% frame in the F-BOR book, CP-19: every command replayed).  Procedures
% TIN-P-01..05; the make runs are replayed by test-pubs-turbo (the printed
% lines are the second run's lines that begin ">>> " / "!!! " or carry a
% pass rate).  TIN-P-04 (make selfhost, K2 == K3) and TIN-P-05 (make ddc)
% are M8 placeholders: DRAFT per D-14 (ADR-0007 DEC-06, DEC-08; FO-2).
% The slogan opens the procedures chapter (operator directive 2026-10-17).
\Gates{test-seed-repro, test-tps-lex-build, test-compiler, test-compiler-os, test-pubs-turbo}
\chapter{The acceptance procedure}
\ix{acceptance procedure}
\IsThisGoodForTheCompany

A Turbo Initech compiler is accepted for release only when every procedure
of this chapter has been carried out, in order, on the release revision of
the sources, and every verification has been met.  Each procedure is
recorded: the build log of each step is kept with the release.

\begin{AcceptanceProcedure}{TIN-P-01}{To confirm that the genesis compiler builds reproducibly}
\PField{Purpose:}The genesis compiler must make the same bytes from the same
source every time.
\PField{Frequency:}For every release candidate.
\PField{Performed by:}The release engineer.
\PField{Prerequisites:}The development system holds the InitechOS sources at
the release revision.
\PField{Materials:}None.
\begin{enumerate}
\item Type the command shown and press \key{Enter}.  The build compiles each
  program of the genesis compiler's test set twice, in two separate
  directories, and compares the results:
\end{enumerate}
\MakeRun{seed-repro}
\PField{Verification:}The last line ends with \texttt{byte-identical across two independent builds}.
\PField{Records retained:}The build log.
\end{AcceptanceProcedure}

\begin{AcceptanceProcedure}{TIN-P-02}{To build the first compiler}
\PField{Purpose:}Make $K_1$, the compiler \file{TPS.COM}, with the genesis
compiler.
\PField{Frequency:}For every release candidate, after TIN-P-01.
\PField{Performed by:}The release engineer.
\PField{Prerequisites:}TIN-P-01 has been verified.
\PField{Materials:}None.
\begin{enumerate}
\item Type the command shown and press \key{Enter}:
\end{enumerate}
\MakeRun{tps-lex-build}
\PField{Verification:}The build reports \code{green}, and the image ends below
6F000h.
\PField{Records retained:}The build log and \file{TPS.COM}.
\end{AcceptanceProcedure}

\begin{AcceptanceProcedure}{TIN-P-03}{To accept the compiler on the corpus}
\PField{Purpose:}Prove that Turbo Initech compiles every corpus program
(Appendix A) to code that prints exactly what the Free Pascal compiler's
code prints, on the development system and on InitechDOS, and that it
compiles its own source deterministically.
\PField{Frequency:}For every release candidate, after TIN-P-02.
\PField{Performed by:}The release engineer.
\PField{Prerequisites:}TIN-P-02 has been verified.
\PField{Materials:}None.
\begin{enumerate}
\item Type the command shown and press \key{Enter}.  The build runs the
  corpus on the development system, compiles the compiler's own source
  twice, runs \file{TPS.COM} on InitechDOS, and finally runs every
  program Turbo Initech compiled in one session of InitechDOS:
\end{enumerate}
\MakeRun{compiler-os}
\PField{Verification:}\texttt{compiler\_host\_pass\_rate=100\% (17/17)} and
\texttt{differential\_pass\_rate=100\% (17/17)} are printed, and the last line
reports \code{GREEN}.
\PField{Records retained:}The build log.
\end{AcceptanceProcedure}

\begin{AcceptanceProcedure}{TIN-P-04}{To produce the self-compilation certificate}
\draft{make selfhost, the K2 == K3 certificate, is an M8 placeholder
(ADR-0007 DEC-06, DEC-08; bead initech-ls4): it needs Turbo Initech
resident on InitechOS with its assembler and linker}
\PField{Purpose:}Prove that the Turbo Initech compiler reproduces itself:
$K_2 = K_3$.
\PField{Frequency:}For every release candidate, after TIN-P-03.
\PField{Performed by:}The release engineer.
\PField{Prerequisites:}TIN-P-03 has been verified.
\PField{Materials:}$K_1$, \file{TPS.COM}, from TIN-P-02.
\begin{enumerate}
\item On InitechOS, compile the Turbo Initech source with $K_1$ and make the
  program file $K_2$.
\item Compile the source with $K_2$ and make $K_3$.
\item Compare $K_2$ and $K_3$ byte for byte.
\end{enumerate}
\PField{Verification:}$K_2$ and $K_3$ are identical.  A difference is not
corrected: it is reported, with both files, as a defect of the compiler or
of the reproducibility of the build.
\PField{Records retained:}$K_2$, $K_3$ and the comparison.
\end{AcceptanceProcedure}

\begin{AcceptanceProcedure}{TIN-P-05}{To compare with an independent genesis compiler}
\draft{make ddc is an M8 placeholder (ADR-0007 DEC-06, FO-2, OQ-2): the
independent genesis compiler does not exist}
\PField{Purpose:}Rule out a fault that reproduces itself through the chain
(Chapter 5).
\PField{Frequency:}For every release.
\PField{Performed by:}The release engineer.
\PField{Prerequisites:}TIN-P-04 has been verified.
\PField{Materials:}The independent genesis compiler.
\begin{enumerate}
\item Make a first compiler with the independent genesis compiler, and
  follow the chain of TIN-P-04 from it.
\item Compare the resulting compiler with $K_3$ of TIN-P-04.
\end{enumerate}
\PField{Verification:}The two compilers are identical.
\PField{Records retained:}Both compilers and the comparison.
\end{AcceptanceProcedure}
